\chapter{jax-Ka}
\label{ch:jax-ka}

% In this chapter we introduce jax-Ka, a differentiable pKa predictor that reads only the protein backbone. Chapter 2 (§2.5) supplies the acid--base background and the structural determinants of pKa shifts; the computational titration models are introduced here, in the chapter preface, because they are the reference points against which jax-Ka is trained and assessed. These are empirical models (PROPKA), Poisson--Boltzmann with Monte Carlo sampling of protonation states (PypKa), constant-pH molecular dynamics, and machine-learned predictors (pKAI and pKAI+, KaML).
% The model is the orientation Graph Query Transformer (oGQT) with relational edges. Residues are nodes of a spatial graph (C-alpha within 20 Å; distance, local-frame direction and relative orientation on each edge; backbone dihedrals on each node). Explicit edges to the residues at sequence offsets of one and two are added, and attention uses a separate key and value projection for each of six relations (i-2, i-1, i+1, i+2, same chain, other chain). Each titratable site is a query token that reads this graph and predicts a bounded shift from its group's reference pKa.
% Training combines pKPDB (PypKa pKas across the PDB) with PINDER protein--protein complexes labelled by pKAI in the bound and free states, so the loss covers absolute pKa, the bound--free shift on binding and an auxiliary burial target. The default frozen checkpoint is the 50% relational model (seed 17, epoch 18); across three seeds the relational edges lowered the selection error from 0.649 to 0.632. An auxiliary target for experimental (crystallisation) pH gave no gain.
% Because the readout differentiates through backbone coordinates, a pKa or a binding-induced pKa shift can be used as a forward model in the same way as the HDX and SAXS models of Chapter 5. MoPrP at pH 4 and pH 7 is the link. The BioFeaturisers simulations need protonation states at each pH, and jax-Ka asks which titratable sites change state between the two conditions and whether those sites lie in the regions that unfold.

% Titration models (preface) → pKMOD and binding-shift state definitions → datasets and masks (pKPDB, PINDER, PKAD) → relational oGQT → binding readout → MoPrP pH 4 / pH 7.

% Superseded plan (pre-October 2026): siamese pKAI retraining and all-atom against backbone Kd-ratio errors (<9%, 30%, <20% over pH 6.1--7.4). Those numbers are not in the current jax-Ka records; do not reuse them without a result to point at.
